\chapter{Recommendations for experimental validation}\label{ch:recommendations}

The following priorities are proposals; none has been carried out. They are ordered so that
each removes an uncertainty on which the next depends.

\paragraph{Priority 1 --- close the chemistry gap.} Characterise the \emph{original}
dichloromethane seed and leaf extracts by LC-MS/MS and GC-MS against authentic standards,
using the published validated UPLC-MS/MS method for niazimicin \parencite{abdelsayed2021jfda}
as a positive control. Until this is done, no compound-level interpretation of the original
result is possible. Yields and concentrations should be reported as measured, never inferred.

\paragraph{Priority 2 --- resolve identities.} Confirm the structures and public identifiers of
the shortlisted compounds, settle the glucomoringin stereochemistry, obtain the paywalled
full texts and recover the missing radiation doses and concentrations. Re-verify the
retraction status of the flagged journal and read the outstanding erratum.

\paragraph{Priority 3 --- complete the deferred statistical validation.} Run the coverage
simulation to choose and justify an interval method, re-assess the primary comparison with
the adopted method, and run at least 1,000 draws of each negative control (permuted outcome
and random gene sets). Only then can superiority be assessed, gene-set attribution be
permitted or any pathway exceed tier~C. The infrastructure for all of this is in place.

\paragraph{Priority 4 --- test purified compounds against whole extracts on the right
endpoint.} Use multi-dose \textbf{clonogenic survival} assays (not 72-hour biomass) with and
without compound or extract, yielding linear-quadratic parameters,
\begin{equation}
  \mathrm{SF}(D) = \exp\bigl(-\alpha D-\beta D^{2}\bigr),
  \label{eq:lq}
\end{equation}
where $\alpha$ and $\beta$ here denote the linear and quadratic radiobiological coefficients,
and a dose-modifying factor for each treatment. Include vehicle-only and cell-free
interference controls, which the original study did not describe, and independent biological
repeats on separate days and passages, since the original replicates were technical wells and
dose was confounded with plate.

\paragraph{Priority 5 --- make the normal-versus-tumour comparison interpretable.} Use
species-matched and tissue-matched normal and tumour pairs assayed identically, so that
species, tissue and malignancy are not confounded as they were when one hamster fibroblast
line was compared with two human carcinoma lines. A selectivity claim requires treated normal
\emph{and} treated tumour measurements on the same endpoint. The best published model of how
such a claim should be worded is the quercetin study of \textcite{horton2013}, which reports
normal-tissue anti-fibrotic protection with \emph{no} effect on fibroblast or murine tumour
radioresponse.

\paragraph{Priority 6 --- test the prioritised pathways directly, if Priority 3 succeeds.}
Perturb the implicated processes (for example p53-pathway activity) in a small panel of cell
lines spanning the predicted range, and measure clonogenic survival. This is hypothesis
testing of a model-derived conjecture, not confirmation.

\paragraph{Priority 7 --- strengthen the benchmark.} Seek a second, independent
radiosensitivity source for cross-assay evaluation, reported as heterogeneity rather than
pooled. Consider prospective prediction on unassayed lines under a panel-selection rule
written into the protocol \emph{before} predictions are viewed, reported as a designed
validation panel rather than an unbiased external cohort.
